\documentclass{article}
\usepackage{caisc_2026}
\usepackage[utf8]{inputenc}
\usepackage[T1]{fontenc}
\usepackage{hyperref}
\usepackage{url}
\usepackage{booktabs}
\usepackage{amsfonts}
\usepackage{nicefrac}
\usepackage{microtype}
\usepackage{xcolor}
\usepackage{graphicx}
\usepackage{multirow}
\usepackage{amsmath}

\title{Inherited vs.\ Intrinsic Irrationality: Do LLMs Exhibit\\
Genuine Behavioral Anomalies or Memorised Human Biases?}

\author{Anonymous Authors}

\begin{document}

\maketitle

\begin{abstract}
Large language models (LLMs) are increasingly deployed as economic agents
in decision-support, financial advisory, and policy simulation
systems---yet it remains unclear whether their apparent behavioral biases
reflect genuine properties of their decision architecture or mere retrieval
of memorised human experimental results. We address this confound
systematically by constructing \emph{isomorphic scenario variants}: novel-domain
and fully abstract reformulations that are mathematically equivalent to five
canonical behavioral economics paradigms (framing effect, loss aversion,
endowment effect, certainty effect, status quo bias) but absent from
standard behavioral economics training corpora. We run controlled
experiments across two architecturally distinct models---GPT-4o-mini (proprietary frontier API)
and Gemma3:4b (open-source 4B-parameter local model)---totalling 15,000 trials across 75 scenario variants at multiple temperatures. Contamination
pre-testing confirms both models accurately describe known human behavioral
patterns prior to experiment exposure (3/3 on all five biases), establishing
memorisation as a genuine confound. Our results reveal three distinct
patterns: (i) some biases collapse on isomorphs, indicating memorisation
(e.g., the open-source model shows zero framing effect on canonical
scenarios, overriding genuine frame-sensitivity observed on novel variants);
(ii) loss aversion is structurally absent across \emph{all} models,
temperatures, and surface types ($\hat{\lambda} \leq 1.05$ vs.\ human
$\lambda \approx 2.25$--$2.5$), suggesting an intrinsic architectural
property independent of training data; and (iii) the two models exhibit
\emph{opposite} deterministic strategies on the canonical Allais paradox
(100\% Allais violation vs.\ 0\%, $p < 0.001$), demonstrating that LLMs
form a heterogeneous class of decision-makers irreducible to either
rational agents or human-like agents. These findings have direct
implications for the deployment of LLMs in high-stakes economic and
policy contexts.
\end{abstract}

\section{Introduction}

A growing body of work examines whether large language models (LLMs)
exhibit systematic deviations from expected utility theory analogous to
those documented in humans~[11,22]. If LLMs reproduce human-like biases,
they may amplify existing irrationality when deployed as decision-support
agents; if they do not, they may serve as rational counterweights to human
bias. Either outcome carries significant implications for AI deployment in
finance~[12], clinical decision-making~[21], and policy simulation~[2].

Recent studies have found evidence of multiple behavioral anomalies in LLMs,
including loss aversion~[16], framing effects~[19], anchoring~[9], and
risk-preference inconsistencies~[3]. However, these studies share a critical
methodological limitation: the \textbf{memorisation confound}. Every canonical
scenario used to elicit these biases---the Asian Disease Problem~[22],
Kahneman and Tversky's mixed gambles~[11], the Allais paradox~[1], Thaler's
mug experiments~[20]---appears verbatim in published behavioral economics
literature that forms a substantial fraction of LLM training corpora~[4].
A model that produces human-like framing effects on the Asian Disease Problem
may be retrieving the known experimental result rather than exhibiting a
genuine preference reversal under equivalent novel conditions.

The most comprehensive existing analysis, by Ross et al.~[16], tests 12 LLMs
including GPT-4 on biases including inequality aversion, loss aversion, and
time discounting---finding GPT-4's loss aversion coefficient ($\lambda = 1.09$)
substantially below the human average ($\lambda \approx 2.56$). Binz and
Schulz~[3] evaluate GPT-3 across 25 cognitive science tasks and find
human-like patterns on some but not others. Yet every scenario in both
studies appears in published literature the models were trained on. McCoy
et al.~[14] demonstrate that LLMs can reproduce benchmark task results
directly from training data, making this concern concrete.

\textbf{This paper introduces the first \emph{isomorph construction}
methodology for behavioral economics---structurally identical scenario variants
with verified mathematical equivalence---enabling systematic separation of
memorised from intrinsic biases.} For each canonical paradigm we construct:
for each canonical paradigm we construct two variants preserving the exact
mathematical structure while replacing the surface domain. A \emph{novel
isomorph} uses a neutral domain absent from behavioral economics literature;
an \emph{abstract isomorph} uses purely symbolic content. If a bias appears
on canonical but collapses on isomorphs, it is memorised. If it persists
uniformly, it is intrinsic.

\paragraph{Contributions.}
(1) A validated isomorph construction methodology with analytical equivalence
verification for five behavioral economics paradigms.
(2) The first cross-model, cross-temperature controlled experiment
distinguishing memorised from intrinsic LLM behavioral anomalies.
(3) Three empirical findings---memorised framing effects, structurally
absent loss aversion, and architecture-dependent Allais strategies---
with direct deployment implications.
(4) Full open-source release of all scenarios, code, and 15,000-trial dataset.

\section{Related Work}

\paragraph{Behavioral biases in LLMs.}
Ross et al.~[16] test 12 LLMs on canonical behavioral economics scenarios,
finding widespread but inconsistent bias patterns across models. Binz and
Schulz~[3] evaluate GPT-3 on 25 cognitive science paradigms, finding
human-like patterns on heuristics and biases but not on reinforcement
learning tasks. Talboy and Fuller~[19] document framing effects across
multiple LLMs but find high prompt sensitivity. Jones et al.~[9] find
anchoring effects in GPT-3, while MacMillan-Scott and Musolesi~[13] report
inconsistent risk preferences across elicitation methods. Santurkar et
al.~[18] find that LLMs exhibit opinion profiles that mirror specific
demographic groups in their training data, raising broader concerns about
memorisation of human behavioral patterns. None of these works control for
training data contamination.

\paragraph{Benchmark contamination and memorisation.}
Carlini et al.~[4] demonstrate that LLMs memorise and can reproduce verbatim
passages from training data at scale. McCoy et al.~[14] show that apparent
reasoning abilities may reflect pattern matching to training examples rather
than genuine competence. Jacovi et al.~[7] propose contamination-aware
evaluation protocols for NLP benchmarks. Deng et al.~[24] document test set
leakage in LLM evaluations and its effect on inflating performance estimates.
Our work extends this contamination literature to behavioral economics,
introducing isomorph construction as a domain-general decontamination method.

\paragraph{Prospect theory and decision theory.}
We ground our analysis in cumulative prospect theory~[23], which parametrises
loss aversion via $\lambda$ (typically $\approx 2.25$--$2.5$ in humans~[23]),
risk sensitivity via $\alpha$ (typically $\approx 0.88$~[11]), and probability
weighting. The Allais paradox~[1] operationalises violations of the
independence axiom. The endowment effect~[20,10] is operationalised via
WTA/WTP divergence, typically $\approx 2\times$ in humans~[10]. Status quo
bias~[17] and default effects~[8] provide our fifth paradigm.

\paragraph{LLMs as economic agents.}
Horton~[6] proposes using LLMs as ``silicon samples'' for economic research,
arguing they can simulate human behavioral patterns. Argyle et al.~[2] find
that LLMs can reproduce survey responses of specific demographic groups
when given appropriate personas. Our findings complicate this picture:
behavioral signatures depend critically on surface domain, suggesting that
using LLMs as behavioral proxies without contamination controls may produce
systematically misleading results.

\section{Methodology}

\subsection{Isomorph Construction}

For each of five canonical behavioral economics paradigms, we constructed
three surface variants sharing identical mathematical structure:

\begin{itemize}
    \item \textbf{Canonical}: The standard scenario as published in the
    behavioral economics literature (e.g., the 600-person disease
    problem~[22]).
    \item \textbf{Novel}: A mathematically equivalent scenario in a neutral,
    low-valence, non-technical domain absent from behavioral economics
    textbooks (e.g., 600 items in a damaged warehouse). Novel domains were
    explicitly chosen to avoid CS-adjacent language, medical framing, and
    financial terminology.
    \item \textbf{Abstract}: A fully abstract formulation with zero semantic
    content (e.g., ``600 units of Resource X'', ``Object Q'', ``State A
    vs.\ State B'').
\end{itemize}

Mathematical equivalence was verified analytically for all 15 scenario
variants prior to data collection. For the framing effect,
$\text{EV}(A) = \text{EV}(B) = 200$ in both gain and loss frames with
mirrored probability ordering. For loss aversion, all gain levels yield
positive EV: $\text{EV}(\text{gamble}) = 0.5G - 50 > 0$ for all
$G \in \{105, 120, 150, 175, 200, 250\}$, making rational acceptance
correct in every trial. For the certainty effect, Choice~2 equals
Choice~1 scaled by $\nicefrac{1}{4}$ (0.80$\to$0.20, 1.00$\to$0.25),
preserving the Allais structure. For the endowment effect, the exchange
wording is perfectly symmetric. For status quo, both options are
explicitly stated to have identical outcomes in all conditions.

\subsection{Experimental Controls}

Four design choices prevent confounds beyond the memorisation confound:
\textbf{position bias control} via prompt-order counterbalancing for loss
aversion (gamble-first vs.\ status-quo-first) and a $2\times2$ orthogonal
design for status quo (default present/absent $\times$ ``stay'' listed
first/second); \textbf{parameter sweeps} for loss aversion
($G \in \{105, 120, 150, 175, 200, 250\}$, loss fixed at 100) and
endowment (\$2, \$5, \$10); \textbf{output format enforcement} via
explicit instructions ending every prompt (parse failure rate across
all 15,000 trials: 0\%); and \textbf{temperature sweep} at
$T \in \{0.0, 0.3, 0.7\}$.

\subsection{Contamination Pre-Testing}

Prior to the main experiment, each model was queried in a \emph{separate
session} with contamination prompts asking it to describe known human
behavioral patterns for each canonical scenario without answering the
scenario itself. Responses were scored 0--3 (0 = no knowledge;
3 = accurate recall of quantitative behavioral statistics). Both models
scored 3/3 on all five biases, confirming that \emph{memorisation risk
is a genuine confound}: both models demonstrably possess the knowledge
required to retrieve canonical results. This pre-test is a necessary but
not sufficient condition for establishing that memorisation occurred in any
given experiment; the primary evidence for memorisation is the
canonical-vs-novel dissociation observed in the experimental results.

\subsection{Models and Procedure}

We tested two models: \textbf{GPT-4o-mini} (proprietary frontier-scale API
model, $T \in \{0.0, 0.7\}$) and \textbf{Gemma3:4b} (open-source
4B-parameter quantised local model, $T \in \{0.3, 0.7\}$). Each of
the 75 scenario variants was presented for 50 independent trials per
temperature condition, with a fresh context window per trial.
Total observations: 15,000.

\subsection{Statistical Analysis}

Binary choice outcomes are compared using $\chi^2$ tests with continuity
correction; effect sizes are Cohen's $h$~[5]. The loss aversion coefficient
$\lambda$ is estimated as the interpolated gain threshold where accept
rate crosses 0.5, divided by the fixed loss of 100; we report upper bounds
given universal acceptance. For the Allais paradox we report the full four-pattern distribution. Significance thresholds: $\alpha = 0.05$ (*), $0.01$ (**), $0.001$ (***).
For key proportions we report 95\% Wilson score confidence intervals where
the bound materially affects interpretation; full CIs for all 136 comparisons
are provided in the supplementary \texttt{stats\_summary.csv}.

\section{Results}

Table~\ref{tab:main} summarises all key findings across models,
surfaces, and temperature conditions.

\begin{table}[t]
\centering
\caption{Summary of key experimental results. GPT-4o-mini = proprietary
frontier API model; Gemma3:4b = open-source 4B local model. All results
at $T=0.7$ unless noted. Significance: *** $p<0.001$, * $p<0.05$,
ns = not significant.}
\label{tab:main}
\small
\begin{tabular}{llcccl}
\toprule
\textbf{Bias} & \textbf{Surface} &
\textbf{GPT-4o-mini} & \textbf{Gemma3:4b} &
\textbf{Human} & \textbf{Interpretation} \\
\midrule
\multirow{3}{*}{\shortstack[l]{Framing\\$\Delta P(\text{A})$}}
  & Canonical & $+0.78^{***}$ & $+0.00^{\text{ns}}$ & $\approx+0.50$ & Memorised (Gemma3); genuine (GPT) \\
  & Novel     & $+0.72^{***}$ & $+1.00^{***}$       & ---            & Genuine both models \\
  & Abstract  & $+0.14^{*}$   & $+0.04^{\text{ns}}$ & ---            & Collapses both models \\
\midrule
\multirow{3}{*}{\shortstack[l]{Loss aversion\\$\hat{\lambda}$}}
  & Canonical & $\leq 1.05$ & $\leq 1.05$ & $2.25$--$2.5$ & Absent: intrinsic \\
  & Novel     & $\leq 1.05$ & $\leq 1.05$ & ---           & Absent: intrinsic \\
  & Abstract  & $\leq 1.05$ & $\leq 1.05$ & ---           & Absent: intrinsic \\
\midrule
\multirow{3}{*}{\shortstack[l]{Endowment\\Owner accept}}
  & Canonical & $0.56^{***}$  & $0.01^{\text{ns}}$ & WTA/WTP $\approx 2\times$ & GPT partial; Gemma3 floor \\
  & Novel     & $0.07^{\text{ns}}$ & $0.00^{\text{ns}}$ & --- & Absent both \\
  & Abstract  & $1.00^{***}$  & $1.00^{***}$        & --- & Ownership signal: both \\
\midrule
\multirow{3}{*}{\shortstack[l]{Allais\\Violation rate}}
  & Canonical & $0.00^{***}$ & $1.00^{***}$ & $\approx 0.65$ & \textbf{Opposite strategies} \\
  & Novel     & $0.36^{***}$ & $0.00^{***}$ & ---            & Dissociation \\
  & Abstract  & $0.18^{*}$   & $0.48^{***}$ & ---            & Mixed \\
\midrule
\multirow{3}{*}{\shortstack[l]{Status quo\\Default effect}}
  & Canonical & $+0.50$ & $+0.50$ & $\approx+0.40$ & Confounded w/ position \\
  & Novel     & $+0.50$ & $+0.50$ & ---            & Confounded w/ position \\
  & Abstract  & $+0.00$ & $+0.00$ & ---            & Absent both models \\
\bottomrule
\end{tabular}
\end{table}

\subsection{Framing Effect}

GPT-4o-mini exhibited a strong canonical framing effect ($\Delta P(\text{A})
= +0.78$, $p < 0.001$, $h = 1.93$) approaching the human baseline
($\Delta \approx +0.50$), and a comparable novel effect ($\Delta = +0.72$,
$p < 0.001$, $h = 1.79$). At $T=0.0$, effects sharpened: canonical
$\Delta = +1.00$ ($h = 3.14$), novel $\Delta = +0.96$ ($h = 2.74$),
suggesting lower temperature reveals the model's deterministic preference
structure. Abstract scenarios yielded a marginal effect at $T=0.7$
($\Delta = +0.14$, $p = 0.036$) that disappeared at $T=0.0$ ($p = 0.24$).

Gemma3:4b showed a canonical-vs-novel dissociation that constitutes the
primary evidence for memorisation. Two components jointly establish this:
(i)~the contamination pre-test confirms Gemma3:4b can accurately describe
the canonical human framing response, establishing knowledge possession;
(ii)~Gemma3:4b shows zero framing effect on canonical scenarios
($\Delta = 0.00$, $p = 1.00$) at both $T=0.3$ and $T=0.7$---choosing
Option~A regardless of frame---while showing a perfect framing effect on
novel variants. The combination of knowledge possession and canonical
suppression is most parsimoniously explained by retrieval of the memorised
``safe choice'' response. On novel scenarios: a
perfect framing effect ($\Delta = +1.00$, $p < 0.001$, $h = 3.14$), showing
Gemma3:4b possesses genuine frame-sensitivity that is overridden by a
memorised response on the canonical scenario. Both models showed zero
framing effect on abstract scenarios at all temperatures.

All framing patterns were fully stable across tested temperatures for both
models.

\begin{figure}[t]
  \centering
  \includegraphics[width=\linewidth]{fig1_framing.png}
  \caption{Framing effect: $P(\text{Choose A})$ under gain frame (solid)
  and loss frame (hatched) for both models across all three surfaces.
  Red dashed/dotted lines indicate human baselines (gain: 0.72, loss:
  0.22~[22]). Significance stars indicate $\chi^2$ test results
  (*** $p<0.001$, * $p<0.05$, ns: not significant).}
  \label{fig:framing}
\end{figure}

\subsection{Loss Aversion}

Both models accepted every gamble at every gain level, across all surfaces
and all temperatures, yielding $\hat{\lambda} \leq 1.05$ in all 24
model-surface combinations (Figure~\ref{fig:lossav}). Both models accepted
a 50\% chance of losing 100 units for a 50\% chance of gaining 105
units---an EV advantage of only 2.5 units---which the vast majority of
human participants would reject under standard prospect theory
parameterisations~[23].

Accept rates were statistically identical under gamble-first and
status-quo-first prompt orderings (difference $< 0.02$ in all conditions),
ruling out position bias. The absence was uniform across canonical, novel,
and abstract surfaces, confirming this is not a memorisation artefact:
it would require memorising a result \emph{contradicting} the published
human literature in the training data.

Comparing to Ross et al.~[16] who report $\lambda = 1.09$ for GPT-4 using
a single canonical gamble pair: our six-point gain sweep across three
surfaces confirms and extends this finding. Not only is $\hat{\lambda}$ low,
it is consistently low regardless of domain, gain magnitude, and prompt
order. \textbf{Loss aversion is structurally absent from both tested
models, independent of training data, domain, temperature, and scale.}
The lowest gain level tested ($G=105$, EV advantage of 2.5 units)
was accepted universally; extending the sweep to lower ratios (e.g.$G=101$) would tighten the bound further but was not necessary to
establish the substantive finding: both models accept gambles that
97.5\% of human participants would reject~[23].

\begin{figure}[t]
  \centering
  \includegraphics[width=0.85\linewidth]{fig2_loss_aversion.png}
  \caption{Loss aversion: estimated $\hat{\lambda}$ upper bounds across
  models and surfaces. All conditions fall at or below $\hat{\lambda}=1.05$
  (grey dashed line: rational agent $\lambda=1.0$; red dashed line: human
  baseline $\lambda\approx 2.25$~[23]).}
  \label{fig:lossav}
\end{figure}

\subsection{Endowment Effect}

\textbf{Abstract surfaces} produced a striking, replicable result: owners
accepted all exchanges (rate = 1.00) while non-owners refused all exchanges
(rate = 0.00), across all three prices and both temperatures for both models
($p < 0.001$ in all cases). This perfect ownership-driven split replicates
across two architecturally distinct models, suggesting that both models respond to stated ownership in abstract contexts, though the underlying mechanisms may differ across architectures.

\textbf{Canonical surfaces} showed divergent model behaviour. GPT-4o-mini showed
a significant price-dependent ownership signal: at \$2, owner accept = 0.67
vs.\ non-owner = 0.00 ($p < 0.001$, $h = 1.82$); at \$10, owner accept
= 1.00 vs.\ non-owner = 0.13 ($p < 0.001$, $h = 2.94$). At $T=0.0$, the ownership ratio increased to 16.7$\times$ at \$10.
The non-monotonic acceptance pattern for GPT-4o-mini owners (0.67 at \$2,
0.33 at \$5, 1.00 at \$10) may reflect threshold reasoning rather than
continuous valuation---the model may treat \$5 as below some internal
reference price for a mug while treating \$10 as meeting it. Gemma3:4b
showed near-floor rates for both owner and non-owner conditions ($< 0.02$),
consistent with treating any price in the \$2--\$10 range as insufficient
for a real object regardless of ownership.

\textbf{Novel surfaces} yielded null effects for both models (all rates
$< 0.10$), suggesting either floor effects or genuine absence of ownership
sensitivity in unfamiliar domains without memorised value associations.

\begin{figure}[t]
  \centering
  \includegraphics[width=\linewidth]{fig3_endowment.png}
  \caption{Endowment effect: acceptance rate curves by ownership condition
  and price for GPT-4o-mini (top row) and Gemma3:4b (bottom row) across all three
  surfaces. Solid lines: owner condition; dashed lines: non-owner condition.
  Human WTA/WTP ratio $\approx 2\times$~[10].}
  \label{fig:endow}
\end{figure}

\subsection{Certainty Effect (Allais Paradox)}

The most striking finding in our dataset involves the two models producing
opposite deterministic strategies on canonical Allais scenarios, stable
across all temperatures (Figure~\ref{fig:certainty}).

\textbf{Gemma3:4b}: 100\% Allais violation (pattern B,A) at $T=0.3$ and
$T=0.7$---choosing the certain option in Choice~1 but the risky option
in Choice~2, \emph{exactly} matching the canonical human violation
documented in the literature~[1,11]. This pattern persisted with zero
variance at both temperatures, indicating a deterministic memorised response.

\textbf{GPT-4o-mini}: 0\% Allais violation, 100\% consistent risk aversion
(pattern B,B) at both $T=0.0$ and $T=0.7$---choosing the lower-variance
option in both choices. This is internally consistent but suboptimal under
expected utility maximisation.

The chi-square test confirms the cross-model difference is highly significant
($p < 0.001$). Neither model produced the rational response (A,A) on any
canonical trial.

On \textbf{novel scenarios}, the pattern shifted: Gemma3:4b switched to 100\%
risk aversion (B,B), providing further evidence that the canonical 100\%
violation reflects memorisation---the model knows the ``correct'' answer for
the canonical Allais problem but defaults to risk aversion on novel variants.
GPT-4o-mini showed a mixed distribution at $T=0.7$ (36\% violation, 62\% risk
aversion) that concentrated toward risk aversion at $T=0.0$ (16\% and 84\%).

On \textbf{abstract scenarios}, GPT-4o-mini at $T=0.0$ produced 100\% reverse
violation (A,B)---a pattern that is internally inconsistent and irrational
under any normative framework---suggesting that low-temperature decoding
exposes a distinct latent bias on fully symbolic inputs.

\begin{figure}[t]
  \centering
  \includegraphics[width=\linewidth]{fig4_certainty.png}
  \caption{Certainty effect: choice pattern distributions across models and
  surfaces. Stacked bars show proportions of each response pattern. Red
  dashed line: human Allais violation rate ($\approx 0.65$~[1]).}
  \label{fig:certainty}
\end{figure}

\subsection{Status Quo Bias}

The $2\times2$ orthogonal design decomposed the apparent status quo effect
into a genuine default effect and a position preference. For both models
on canonical and novel surfaces: stay rates were 1.00 in Conditions
C1 (default + stay first), C2 (default + stay second), and C3
(no default + stay first), and 0.00 in Condition C4 (no default +
switch first). The decomposed default effect (+0.50) and position effect
(+0.50) are equal---the only condition producing non-stay responses is C4,
where ``switch'' appears as Option~A. This demonstrates that both models
exhibit a strong, domain-general preference for Option~A that is
independent of whether any default has been set.

On abstract surfaces, GPT-4o-mini showed pure position bias (position
effect $= +1.00$, default effect $= 0.00$), confirming the Option~A
preference extends to fully symbolic contexts. Gemma3:4b showed zero
effects in either direction on abstract scenarios.

This finding is itself informative: our $2\times2$ design cannot isolate a
genuine default effect from the dominant Option~A preference---these two
signals are statistically inseparable in our data. Whether LLMs lack the
mechanism underlying human default effects~[8,17], or whether that mechanism
is simply swamped by a stronger position preference, cannot be determined from
the current design. What is clear is that standard two-option experimental
designs without position counterbalancing will conflate position preference
with status quo bias when testing LLMs.

\begin{figure}[t]
  \centering
  \includegraphics[width=\linewidth]{fig5_status_quo.png}
  \caption{Status quo bias: heatmaps of default effect (left) and position
  effect (right) across both models and all three surfaces. Equal default
  and position effects of $+0.50$ on canonical and novel surfaces indicate
  the two cannot be separated, revealing dominant Option~A preference rather
  than genuine default-driven bias.}
  \label{fig:statusquo}
\end{figure}

\section{Discussion}

\subsection{A Tripartite Taxonomy of LLM Behavioral Anomalies}

Our results support a taxonomy of three qualitatively distinct patterns.
Figure~\ref{fig:summary} provides a cross-model overview of all key metrics
relative to human baselines.

\begin{figure}[t]
  \centering
  \includegraphics[width=\linewidth]{fig6_summary_heatmap.png}
  \caption{Cross-model summary heatmap: key metrics for both models relative
  to human baselines. Cell values show raw metric; parenthetical shows human
  baseline. Colour intensity reflects proximity to human baseline. Loss aversion
  $\hat{\lambda}$ is consistently far below human levels across all conditions.}
  \label{fig:summary}
\end{figure}

\paragraph{Memorised biases.}
These appear on canonical scenarios but collapse on isomorphs. Gemma3:4b's
framing effect dissociation (zero effect canonical, perfect effect novel)
and its canonical Allais 100\% violation (absent on novel) both exemplify
this pattern. The contamination scores (3/3 on all five biases for both
models) confirm the mechanism: models recognise the canonical scenario and
retrieve the known result, overriding whatever reasoning they apply to
novel equivalents. This finding suggests that existing reports of LLM
behavioral biases using only canonical scenarios~[16,3,19] may be
substantially overstating the prevalence of genuine bias.

\paragraph{Intrinsic structural properties.}
Loss aversion is the unambiguous example: $\hat{\lambda} \leq 1.05$
across all 24 model-surface-temperature combinations. This cannot be a
memorisation artefact because it contradicts the human-calibrated training
data. We interpret this as reflecting how these models process probabilistic
outcomes: in the absence of affective weighting or evolved risk sensitivity,
they approximate expected value maximisation~[3], accepting any positive-EV
gamble regardless of loss magnitude. The abstract-surface endowment result
(ownership separation at 100\%/0\%) may be a second example of an intrinsic
property, though more investigation is needed.

\paragraph{Architecture-dependent signatures.}
The canonical Allais reversal---100\% violation vs.\ 0\%---cannot both
be memorisation of the same training data. We hypothesise that differences
in instruction tuning objectives, RLHF methodology~[15], and model scale
produce qualitatively different decision heuristics on problems with
certainty effects. GPT-4o-mini's 100\% reverse violation (A,B) at $T=0.0$ on
abstract scenarios suggests an additional latent bias that only manifests
under deterministic decoding, warranting further investigation.

\subsection{The Option~A Position Bias}

Both models exhibit a strong, domain-general preference for Option~A
that is temperature-stable and independent of whether a default has been
set. This can masquerade as status quo bias in designs that do not
counterbalance position. The implication is methodological: any study
presenting LLMs with two-option choices labelled A and B without position
counterbalancing risks confounding genuine preference signals with this
structural labelling bias~[13].

\subsection{Implications for LLM Deployment}

The structurally absent loss aversion has the most direct practical
implications. LLM-assisted systems in portfolio optimisation, insurance
pricing, or clinical treatment choice cannot assume the model will
weight losses more heavily than equivalent gains. Under prospect theory
parameterisations calibrated to human data~[23], both models would
systematically accept gambles that human clients would reject.

The framing result implies that GPT-4o-mini's genuine frame-sensitivity
on novel scenarios may reasonably approximate human responses, while
Gemma3:4b reproduces the canonical human framing effect for memorisation
reasons that do not generalise to novel policy scenarios. Using these
models interchangeably as behavioral proxies~[6,2] would produce
contradictory predictions on novel contexts.

\subsection{Limitations}

We note four limitations bound the scope of our conclusions.
\emph{Model coverage}: we tested two models representing API-scale and
4B local deployment; findings cannot be assumed to generalise to other
model families, scales, or training regimes without further testing.
\emph{Endowment price sensitivity}: the \$2--\$10 sweep may fall below
the psychological salience threshold for real-world objects in both models,
producing floor effects that prevent measurement of the endowment effect
rather than evidencing its absence.
\emph{Single-turn prompting}: all trials used zero-context single-turn
prompts without role assignment or system-prompt priming; behavioral
signatures may differ under persona-assigned or multi-turn conditions
that more closely replicate real deployment contexts~[2,6].
\emph{Certainty effect domain coverage}: the Allais scenarios were tested
with monetary framing only; whether the architecture-dependent violation
patterns extend to non-monetary contexts remains an open question.

\section{Conclusion}

We introduced and validated an isomorph-based methodology for distinguishing
memorised from intrinsic behavioral anomalies in LLMs, and applied it across
five canonical behavioral economics paradigms, two architecturally distinct
models, and three temperature conditions totalling 15,000 trials. LLMs are
neither rational agents nor human-like agents: they exhibit a tripartite
taxonomy comprising memorised biases, intrinsic structural properties, and
architecture-dependent signatures. Loss aversion is structurally absent from
both tested models across all conditions---our most robust finding, with
immediate implications for any deployment context assuming human-calibrated
risk sensitivity. The canonical Allais reversal shows that behavioral
characterisation cannot be generalised across LLM families. The Option~A
position bias reveals a methodological confound in standard two-option
experimental designs. We release all code, scenarios, and data to support
future work on behaviorally-grounded LLM evaluation.

\section*{References}

\small

[1] Allais, M. (1953). Le comportement de l'homme rationnel devant le risque.
\textit{Econometrica}, 21(4), 503--546.

[2] Argyle, L.P., et al. (2023). Out of one, many: Using language models to
simulate human samples. \textit{Political Analysis}, 31(3), 337--351.

[3] Binz, M., \& Schulz, E. (2023). Using cognitive psychology to understand
GPT-3. \textit{PNAS}, 120(6), e2218523120.

[4] Carlini, N., et al. (2021). Extracting training data from large language
models. \textit{USENIX Security Symposium}, 2633--2650.

[5] Cohen, J. (1988). \textit{Statistical power analysis for the behavioral
sciences} (2nd ed.). Lawrence Erlbaum.

[6] Horton, J.J. (2023). Large language models as simulated economic agents:
What can we learn from homo silicus? \textit{NBER Working Paper} 31122.

[7] Jacovi, A., et al. (2023). Stop uploading test data in plain text:
Practical strategies for mitigating data contamination. \textit{EMNLP 2023}.

[8] Johnson, E.J., \& Goldstein, D. (2003). Do defaults save lives?
\textit{Science}, 302(5649), 1338--1339.

[9] Jones, E., \& Steinhardt, J. (2022). Capturing failures of large language models
via human cognitive biases. \textit{Advances in Neural Information Processing Systems},
35, 11785--11799. arXiv:2202.12299.

[10] Kahneman, D., Knetsch, J.L., \& Thaler, R.H. (1991). Anomalies: The
endowment effect, loss aversion, and status quo bias. \textit{Journal of
Economic Perspectives}, 5(1), 193--206.

[11] Kahneman, D., \& Tversky, A. (1979). Prospect theory: An analysis of
decision under risk. \textit{Econometrica}, 47(2), 263--291.

[12] L\'{o}pez-Lira, T., \& Tang, Y. (2023). Can ChatGPT forecast stock
price movements? \textit{arXiv:2304.07619}.

[13] MacMillan-Scott, O., \& Musolesi, M. (2024). (Ir)rationality and
cognitive biases in large language models. \textit{Royal Society Open
Science}, 11, 240255.

[14] McCoy, R.T., et al. (2023). Embers of autoregression: Understanding
large language models through the problem they are trained to solve.
\textit{arXiv:2309.13638}.

[15] Ouyang, L., et al. (2022). Training language models to follow
instructions with human feedback. \textit{NeurIPS 2022}.

[16] Ross, J., Kim, Y., \& Lo, A.W. (2024). LLM economicus? Mapping the behavioral
biases of LLMs via utility theory. \textit{Conference on Language Modeling (COLM 2024)}.
arXiv:2408.02784.

[17] Samuelson, W., \& Zeckhauser, R. (1988). Status quo bias in decision
making. \textit{Journal of Risk and Uncertainty}, 1(1), 7--59.

[18] Santurkar, S., et al. (2023). Whose opinions do language models
reflect? \textit{ICML 2023}.

[19] Talboy, A.N., \& Fuller, E. (2023). Challenging the appearance of
machine intelligence: Cognitive bias in LLMs. \textit{arXiv:2304.13734}.

[20] Thaler, R. (1980). Toward a positive theory of consumer choice.
\textit{Journal of Economic Behavior \& Organization}, 1(1), 39--60.

[21] Thirunavukarasu, A.J., et al. (2023). Large language models in
medicine. \textit{Nature Medicine}, 29, 1930--1940.

[22] Tversky, A., \& Kahneman, D. (1981). The framing of decisions and
the psychology of choice. \textit{Science}, 211(4481), 453--458.

[23] Tversky, A., \& Kahneman, D. (1992). Advances in prospect theory:
Cumulative representation of uncertainty. \textit{Journal of Risk and
Uncertainty}, 5(4), 297--323.

[24] Deng, C., Zhao, Y., Tang, X., Gerstein, M., \& Cohan, A. (2023). Investigating
data contamination in modern benchmarks for large language models. arXiv:2311.09783.

%%%%%%%%%%%%%%%%%%%%%%%%%%%%%%%%%%%%%%%%%%%%%%%%%%%%%%%%%%%%
\newpage

\section*{Conference For AI Scientists 2026 - AI Involvement Checklist}
\label{checklist_ai_involvement}

\subsection*{Research Stage Assessment}

\begin{enumerate}

\item \textbf{Hypothesis development}

Answer: \involvementC{}

Iteration Effort: \iterationMedium{}

Explanation: The research question---distinguishing memorised from intrinsic biases via isomorph
construction---was proposed by a frontier LLM acting as a research collaborator, following its own
literature analysis identifying the memorisation confound. The LLM generated the initial three-finding
taxonomy and isomorph construction procedure. The human validated the framing, selected the
five paradigms, and made the final decision to pursue the study. Approximately 20--30 substantive
iterations to sharpen the hypothesis.

\item \textbf{Experimental design and implementation}

Answer: \involvementC{}

Iteration Effort: \iterationHigh{}

Explanation: The LLM designed all 15 isomorph scenario variants through six rounds of reviewer-style
critique and revision. The LLM wrote all experimental Python scripts: scenario generation, API
harnesses for two cloud and one local provider, response parsers, rate-limit handling, and the full
analysis and figure generation pipeline. The human verified mathematical equivalence of all isomorphs,
selected final domain wordings, diagnosed API integration issues, and executed all 15,000 experimental
trials. Over 100 substantive iterations across scenario design, code debugging, and API troubleshooting.

\item \textbf{Analysis of data and interpretation of results}

Answer: \involvementC{}

Iteration Effort: \iterationMedium{}

Explanation: The LLM wrote the full statistical analysis pipeline including $\chi^2$ tests, Cohen's $h$,
prospect theory $\lambda$ estimation, and all six publication figures. The LLM provided initial
interpretations and proposed the tripartite taxonomy. The human reviewed all statistical outputs,
validated the $\lambda$ estimation approach, and refined the Discussion argumentation. Approximately
15--25 substantive iterations.

\item \textbf{Writing}

Answer: \involvementC{}

Iteration Effort: \iterationMedium{}

Explanation: The LLM produced the first complete draft of the paper including all sections, tables,
figures, references, and both mandatory checklists. The human reviewed for scientific accuracy,
verified all cited statistics match experimental outputs, added citations, tightened arguments, and made
all final editorial decisions. Approximately 10--20 substantive iterations.

\end{enumerate}

\subsection*{AI System Documentation}

\begin{enumerate}
\setcounter{enumi}{4}

\item \textbf{AI systems used}

Description: All AI-assisted research stages used a single frontier large language model accessed via a
consumer chat interface. The system is a transformer-based instruction-tuned model at frontier scale
(estimated 100B+ parameters) with broad scientific, coding, and statistical analysis capabilities. The
same system was used for all stages. No multi-agent framework or autonomous pipeline was used; all
interactions were human-directed short-context conversations with outputs approved before use.

\item \textbf{Observed AI Limitations}

Description: Four recurring limitations: (1)~scenario wordings that introduced subtle semantic
confounds (CS-adjacent terminology, asymmetric verb framing) requiring human identification across
multiple review cycles; (2)~a statistical code error in $\lambda$ upper-bound handling for the
degenerate all-accept case, requiring debugging; (3)~overconfident causal claims in early Discussion
drafts requiring moderation; (4)~inability to independently execute API calls or verify experimental
results---all runs required human execution with results uploaded for analysis.

\end{enumerate}

\newpage

\section*{Conference For AI Scientists 2026 - Reproducibility and Responsibility Checklist}
\label{checklist_reproducibility}

\begin{enumerate}

\item \textbf{Claims}
\item[] Answer: \answerYes{}
\item[] Justification: The abstract and introduction accurately state the three core empirical findings.
All quantitative claims in the abstract are verified against experimental output files. Claims are
explicitly scoped to the two tested models throughout the paper.

\item \textbf{Limitations}
\item[] Answer: \answerYes{}
\item[] Justification: Section~5.4 contains a dedicated Limitations paragraph addressing four
substantive limitations: model coverage scope, endowment price sensitivity, single-turn prompting
context, and certainty effect domain coverage.

\item \textbf{Theory assumptions and proofs}
\item[] Answer: \answerNA{}
\item[] Justification: No theoretical claims require formal proof. Mathematical equivalence of all 15
isomorph variants is verified analytically in Section~3.1.

\item \textbf{Experimental result reproducibility}
\item[] Answer: \answerYes{}
\item[] Justification: All prompts, model configurations, temperature settings, trial counts, and
analysis procedures are described in Section~3 and supplementary materials. All raw response data
is released publicly.

\item \textbf{Open access to data and code}
\item[] Answer: \answerYes{}
\item[] Justification: All 15,000 trial responses, 75 scenario definitions, experimental scripts, and
analysis code are released anonymously as supplementary material. The repository README contains
exact commands for reproducing all results.

\item \textbf{Experimental setting/details}
\item[] Answer: \answerYes{}
\item[] Justification: Section~3 specifies all experimental conditions: temperatures, trials per
condition, prompt format instructions, API configurations, and gain/price sweep parameters. Full
scenario text for all 75 variants is in the supplementary dataset.

\item \textbf{Experiment statistical significance}
\item[] Answer: \answerYes{}
\item[] Justification: All key comparisons report $\chi^2$ test $p$-values with significance stars
and Cohen's $h$ effect sizes. The full 136-row statistical summary is provided as
\texttt{stats\_summary.csv} in the supplementary materials.

\item \textbf{Experiments compute resources}
\item[] Answer: \answerYes{}
\item[] Justification: Gemma3:4b was run locally on a consumer GPU (4GB VRAM, 8GB RAM, 3.2GHz)
requiring approximately 8--9 hours per 3,750-trial run. GPT-4o-mini was accessed via cloud API
at approximately \$0.05--\$0.10 per 3,750-trial run with average latency 700--1100ms per call.
Total compute cost was under \$1 USD.

\item \textbf{Research Ethics}
\item[] Answer: \answerYes{}
\item[] Justification: This study involves no human subjects, personal data, or sensitive information.
All experiments were conducted solely on language model APIs and local inference. The work conforms
with the NeurIPS Code of Ethics.

\item \textbf{Broader impacts}
\item[] Answer: \answerYes{}
\item[] Justification: Positive impacts include providing practitioners with empirically grounded
guidance for LLM deployment in economic decision contexts. Potential misuse (exploiting predictable
model biases) is limited by the specificity and context-dependence of the findings. Deployment
implications are discussed in Section~5.3.

\end{enumerate}

\end{document}
